\section{Unified solver mechanics as method-evolution objects}
\label{sec:unified_solver_learning}

The unified problem-solving projection adds several proposal-only method objects---operational-map coverage, path/concurrency quotienting, multiobjective path cost, fieldability, Mechanic Differential Diagnosis (MDD), Verified Solver Compilation (VSC) and trajectory-to-certificate assembly. They do not constitute a second experience system. Their executed decisions, costs, residuals and verification outcomes are recorded through the existing \TaskEpisode{} substrate and may become \Lesson{} or ResearchTool candidates only through the consolidation and fresh-assurance process defined in this paper.

This matters most for MDD and VSC. A diagnosis that labels a stalled route as a representation, metric, scale, verifier or operator-basis problem is merely a hypothesis until a discriminator or downstream repair provides evidence. Likewise, a compiled representation--solver--decoder--verifier configuration is a method candidate, not a capability upgrade. Repeated development success may justify a scoped procedural lesson such as ``this transformation effect reduces search in this structural regime,'' but promotion requires held-out transfer, cost accounting, negative-history preservation and the ordinary protected method-evolution gate.

The learning target should therefore remain vector-valued. Useful telemetry includes cause-identification accuracy, specialist conditional advantage, representation/portal preservation, geometry construction and reuse cost, closed-loop routing success, verifier scheduling cost, staleness and root-level residual reduction. A single end-to-end win cannot retroactively attribute every component. Conversely, a failed geometry or solver compilation becomes structured negative history that may reveal a missing operator, ontology coordinate or context boundary. This makes the unified solver a consumer and producer of method-evolution evidence without making it a sovereign evaluator of its own improvement.
